%! Complex Numbers — Unit 3, Lesson 3.4 — Homework
\documentclass[10pt]{article}
\usepackage{algebra2-article}
\usepackage{algebra2-boxes}
\newcommand{\TallMath}[1]{$\displaystyle #1\rule[-1.4em]{0pt}{3.2em}$}

\begin{document}

\pageheader{Unit 3, Lesson 3.4}{Homework}
\namedateperiod

\vspace{0.10in}

\begin{notesbox}{Practice}
Show your work. Pull the $i$ out \emph{first}, replace $i^2=-1$ when you multiply, and write every answer
in standard $a+bi$ form.

\begin{enumerate}[leftmargin=1.9em, itemsep=10pt]
  \item \textbf{Rewrite with $i$} (simplest radical form).
    \[ \sqrt{-25}=\blank{1.6cm} \quad \sqrt{-8}=\blank{1.8cm} \quad \sqrt{-72}=\blank{1.8cm} \quad \sqrt{-100}=\blank{1.6cm} \quad -\sqrt{-49}=\blank{1.6cm} \]

  \item \textbf{Write in $a+bi$ form.}
    \[ 4+\sqrt{-9}=\blank{2.2cm} \quad -2-\sqrt{-16}=\blank{2.2cm} \quad \sqrt{-121}=\blank{2.0cm} \quad 10=\blank{2.0cm} \]

  \item \textbf{Add or subtract.}
    \begin{enumerate}[label=(\alph*), leftmargin=1.9em, itemsep=5pt]
      \item $(5+3i)+(2-7i)=$ \blank{3.0cm}
      \item $(-4+i)-(3-6i)=$ \blank{3.0cm}
      \item $(8-2i)-(8+2i)=$ \blank{3.0cm}
    \end{enumerate}

  \item \textbf{Multiply.} FOIL, then replace $i^2=-1$.
    \begin{enumerate}[label=(\alph*), leftmargin=1.9em, itemsep=5pt]
      \item $(2+i)(3+4i)=$ \blank{3.0cm}
      \item $(5-3i)(2-i)=$ \blank{3.0cm}
      \item $(1+6i)^2=$ \blank{3.0cm}
      \item $3i(4-5i)=$ \blank{3.0cm}
    \end{enumerate}

  \item \textbf{Conjugate products.} Multiply by the conjugate; your answer should be a real number.
    \[ (7+2i)(7-2i)=\blank{2.4cm} \qquad (3-5i)(3+5i)=\blank{2.4cm} \]

  \item \textbf{Powers of $i$.} Reduce the exponent mod $4$.
    \[ i^{12}=\blank{1.4cm} \quad i^{27}=\blank{1.4cm} \quad i^{50}=\blank{1.4cm} \quad i^{101}=\blank{1.4cm} \]

  \item \textbf{Solve over the complex numbers.} Both solutions, in $a+bi$ form.
    \[ x^2=-64:\ x=\blank{2.6cm} \qquad x^2+20=0:\ x=\blank{2.8cm} \]

  \item \textbf{Explain.} A classmate claims ``$i^2=1$, because squaring always makes a number positive.''
        What is wrong with this reasoning, and what is $i^2$ really? \writeline
\end{enumerate}
\end{notesbox}

\begin{extensionbox}
\textbf{Complex roots by completing the square.} Solve each; the two solutions are complex conjugates.
\begin{enumerate}[label=(\alph*), leftmargin=1.8em, itemsep=5pt]
  \item $x^2-6x+13=0 \Rightarrow (x-3)^2=$ \blank{1.4cm}, so $x=$ \blank{3.0cm}
  \item $2x^2+8=0$ (divide by $2$ first) $\Rightarrow x=$ \blank{3.0cm}
  \item Check (a): substitute one of your roots back into $x^2-6x+13$ and show it gives $0$.
\end{enumerate}
\writelines{2}
\end{extensionbox}

\begin{spiralbox}
\textbf{Coming next --- Lesson 3.5: The Quadratic Formula \& the Discriminant.} You can now handle a
\emph{negative} under a square root, so \emph{every} quadratic is solvable. Next you'll get one formula
that solves them all, $x=\dfrac{-b\pm\sqrt{b^2-4ac}}{2a}$, and use the \textbf{discriminant} $b^2-4ac$ to
predict --- before solving --- whether the roots are two real, one real, or two complex conjugates.
\end{spiralbox}

\end{document}
